\documentclass[11pt]{article}
\usepackage[margin=1in]{geometry}
\usepackage{setspace}
\usepackage{amsmath,amssymb}
\usepackage{booktabs}
\usepackage{hyperref}
\usepackage{enumitem}

\hypersetup{
  colorlinks=true,
  linkcolor=blue,
  citecolor=blue,
  urlcolor=blue
}

\title{Physics-Aware Action-Conditioned World Modeling for Billiards\\
\large DiT + Diffusion Forcing on Large Synthetic Multi-Run Datasets}
\author{Moin Arz Mattar}
\date{April 2026}

\begin{document}
\maketitle
\doublespacing

\begin{abstract}
This project studies whether diffusion transformers can preserve physics in an interactive, action-conditioned billiards simulator, and whether diffusion forcing improves long-horizon rollouts. We build a full pipeline: simulation data, VAE latent compression, DiT world-model training, rollout evaluation, and interactive inference tools. The core question is whether multi-step diffusion forcing preserves game logic better than single-step diffusion training. Across multiple runs and model scales, diffusion forcing strongly improves long-horizon stability. Remaining failure modes are mostly identity consistency (color drift and occasional object swallowing), which motivates targeted follow-up losses.
\end{abstract}

\section{Introduction}
World models predict future observations conditioned on actions, enabling simulation and interactive generation. Recent diffusion-transformer systems can look realistic, but preserving physical logic over rollout is still hard. We focus on billiards because it gives explicit checks: do balls bounce correctly, collide correctly, and remain consistent identities over time?

Our novelty is to treat \textit{interactive physics consistency} as the main target, not only visual quality. We test single-step DiT versus DiT+diffusion forcing under the same data and evaluation pipeline, and we inspect both metric trends and failure modes.

\section{Related Work}
Our approach builds on DDPM \cite{ho2020ddpm}, DDIM \cite{song2021ddim}, latent diffusion \cite{rombach2022ldm}, DiT \cite{peebles2023dit}, and VAE latent modeling \cite{kingma2014vae}. We also draw motivation from world models \cite{ha2018worldmodels, hafner2023dreamerv3} and action-conditioned diffusion gameplay work such as Oasis \cite{oasisrepo, oasisblog}.

\section{Problem Setup}
At each time step $t$, we observe:
\begin{align}
x_t \in \mathbb{R}^{72\times 128\times 3},\quad
a_t \in \mathbb{R}^{3},\quad
s_t \in \mathbb{R}^{16\times 4},
\end{align}
where $a_t=[f_x,f_y,\text{trigger}]$ and $s_t$ stores per-ball simulator state (position and velocity channels).

We train a world model to generate $x_{t+1:t+H}$ conditioned on past context and future actions. In practice we work in latent space $z_t\in\mathbb{R}^{4\times 9\times 16}$ for efficiency.

\section{Data and Pipeline}
\subsection{Data generation}
We generated three dataset lineages with different goals:
\begin{itemize}[leftmargin=*]
  \item \textbf{v1} (100,000 episodes): baseline coverage.
  \item \textbf{v2} (100,000 episodes): more diverse setups and trajectories.
  \item \textbf{v3} (200,000 episodes): collision-heavy distribution for hard-contact dynamics.
\end{itemize}
Each episode has 600 frames, for 400,000 episodes total. This staged dataset design was central to improving robustness.

\subsection{VAE stage}
Frames are compressed by a VAE to 4-channel $9\times 16$ latents (about $48\times$ element compression). We use this latent space for all world-model training and decode for evaluation and visualization.

\subsection{Action-conditioned diffusion model}
Given context length $L$, the model receives $z_{t-L+1:t}$, action $a_t$, and noisy target latent $x_t^{(\tau)}$, and predicts noise $\hat{\epsilon}_\theta$.

Forward diffusion:
\begin{align}
x_\tau = \sqrt{\bar{\alpha}_\tau}\,x_0 + \sqrt{1-\bar{\alpha}_\tau}\,\epsilon,\quad \epsilon\sim\mathcal{N}(0,I).
\end{align}

Single-step baseline loss:
\begin{align}
\mathcal{L}_{\text{base}} = \mathbb{E}\left[\|\hat{\epsilon}_\theta - \epsilon\|_2^2\right].
\end{align}

\subsection{Diffusion forcing}
For rollout length $K$, we train over multiple future steps:
\begin{align}
\mathcal{L}_k = \|\hat{\epsilon}_k-\epsilon_k\|_2^2,\quad
\mathcal{L}_{\text{DF}} = \frac{1}{K}\sum_{k=1}^{K}\mathcal{L}_k.
\end{align}
Predicted clean latent estimate:
\begin{align}
\hat{x}_{0,k} = \frac{x_{\tau,k} - \sqrt{1-\bar{\alpha}_\tau}\,\hat{\epsilon}_k}{\sqrt{\bar{\alpha}_\tau}}.
\end{align}
Context for the next step is updated with scheduled teacher forcing:
\begin{align}
\tilde{z}_{k+1}=m_k z_{k+1}^{\text{gt}}+(1-m_k)\hat{x}_{0,k},\quad
m_k\sim\text{Bernoulli}(p_{\text{tf}}(n)),
\end{align}
where $p_{\text{tf}}(n)$ decays over processed samples $n$. This reduces train-test mismatch by exposing training to self-generated contexts.

\section{Experimental Setup}
\subsection{Model and training}
We trained multiple scales, culminating in a 1.52B-parameter DiT+DF model ($d_{\text{model}}=2048$, 30 layers, 32 heads) on 2x H100 using distributed data parallel training.

\subsection{Evaluation protocol}
We use fixed-horizon rollout eval with latent and frame metrics at horizons
$h\in\{1,4,8,16,32\}$:
\begin{itemize}[leftmargin=*]
  \item latent MSE and latent MAE
  \item frame PSNR and frame L1
\end{itemize}
Inference sampling uses DDIM \cite{song2021ddim} for practical speed/quality tradeoffs.

\section{Results}
\subsection{Baseline vs diffusion forcing}
The clearest result is baseline versus DF at context length 8:
\begin{table}[h]
\centering
\begin{tabular}{lcccc}
\toprule
Model & MSE@16 & MSE@32 & PSNR@16 & PSNR@32 \\
\midrule
DiT baseline & 54.023 & 318.312 & 17.133 & 14.857 \\
DiT + DF & 0.763 & 0.827 & 19.447 & 19.372 \\
\bottomrule
\end{tabular}
\caption{Long-horizon rollout quality (lower MSE better, higher PSNR better).}
\label{tab:baseline_df}
\end{table}

DF improves long-horizon stability by a large margin. Qualitatively, trajectories and collisions remain coherent for longer rollouts.

\subsection{Scaling and checkpoint progression}
Scaling from mid-size to larger models improved short- and mid-horizon quality, though not all long-horizon metrics were monotonic.

\subsection{Interactive physics findings}
In interactive tests, wall-bounce behavior was generally correct and collision behavior improved substantially after DF plus v3 collision-heavy data. The main remaining issue is identity consistency: some balls still drift in color or occasionally swallow/merge under complex contact. This suggests dynamics are learned better than persistent appearance identity.

Qualitative media is included as supplementary GIFs: VAE reconstruction comparison and latest model rollout preview.

\section{Novelty and Main Discovery}
Our novelty is the explicit hypothesis test in a controlled physics game:
\begin{quote}
Can diffusion transformers preserve physical logic better when trained with diffusion forcing rather than single-step diffusion loss?
\end{quote}

Our experiments suggest yes, especially at long horizons where error accumulation dominates. Diffusion forcing trains under conditions closer to inference and improves rollout robustness in interactive physics scenarios.

\section{Limitations and Future Work}
\subsection{Current limitations}
\begin{itemize}[leftmargin=*]
  \item Identity consistency (color/object permanence) remains weaker than motion prediction.
  \item Interactive inference is still slower than true real-time gameplay.
  \item Distillation path is still under stabilization.
\end{itemize}

\subsection{Planned next steps}
\begin{itemize}[leftmargin=*]
  \item Add color-consistency auxiliary loss on predicted $x_0$ (ball-focused, simulator-mask-aware).
  \item Increase collision-focused sampling while keeping uniform sampling mix.
  \item Complete student distillation for lower-latency interactive inference.
  \item Evaluate DDIM-step and horizon tradeoffs for playable experience.
\end{itemize}

\section{Conclusion}
This project demonstrates that DiT+DF is an effective direction for action-conditioned billiards world modeling. The main takeaway is that diffusion forcing materially improves long-horizon physical stability versus single-step training. We reached credible wall-bounce and collision behavior, while identifying remaining bottlenecks in identity consistency and inference speed. The next phase is targeted identity losses and stabilized distillation for faster interactive rollout.

\bibliographystyle{plain}
\bibliography{refs}

\end{document}
